\section{Experimental Results and Discussion}
\label{sec:results}

\subsection{Formal Definitions of Metrics}
\label{sec:metrics}

All timing measurements are reported in milliseconds (ms). For each experimental condition, we collect $n = 100$ measured runs after discarding $n_{\text{warmup}} = 10$ warm-up runs. Let $x_1, \dots, x_n$ denote the execution times of the measured runs. The formal definitions of the median and aggregated metrics are given in Equations~\eqref{eq:median}--\eqref{eq:speedup}.

\begin{equation}
\label{eq:median}
\mathrm{median}(x) =
\begin{cases}
x_{((n+1)/2)}, & n \text{ odd}, \\[4pt]
\bigl(x_{(n/2)} + x_{((n/2)+1)}\bigr)/2, & n \text{ even},
\end{cases}
\end{equation}
where $x_{(1)} \le x_{(2)} \le \dots \le x_{(n)}$ is the ordered sample.

\begin{equation}
\label{eq:percentile}
p_k = x_{(\lceil k \cdot n / 100 \rceil)}.
\end{equation}

\begin{equation}
\label{eq:ci}
CI_{95} = [\,\theta_{0.025},\ \theta_{0.975}\,],
\end{equation}
where $\theta_{0.025}$ and $\theta_{0.975}$ are the 2.5th and 97.5th percentiles of the bootstrap medians obtained from $B = 5000$ resamples with replacement (seed = 42).

\begin{equation}
\label{eq:speedup}
S = \frac{\mathrm{median}(x_{\text{baseline}})}{\mathrm{median}(x_{\text{proposed}})}.
\end{equation}
A value $S > 1$ indicates that the proposed method is faster than the baseline; e.g., $S = 4.26$ means the proposed method is 4.26 times faster.

\subsection{Notation}
\label{sec:notation}

Table~\ref{tab:notation} summarizes the symbols used in Equations~\eqref{eq:median}--\eqref{eq:speedup}.

\begin{table}[htbp]
\centering
\caption{Symbols used in the formal definitions of metrics.}
\label{tab:notation}
\begin{tabular}{lll}
\toprule
Symbol & Description & Type \\
\midrule
$n$ & Number of measured runs per condition & scalar (integer) \\
$n_{\text{warmup}}$ & Number of discarded warm-up runs & scalar (integer) \\
$x_i$ & Execution time of the $i$-th measured run (ms) & scalar \\
$x_{(i)}$ & $i$-th order statistic of the sample & scalar \\
$k$ & Percentile level (25, 75, 95, 99) & scalar \\
$p_k$ & $k$-th percentile of execution time (ms) & scalar \\
$B$ & Number of bootstrap resamples & scalar (integer) \\
$CI_{95}$ & Bootstrap 95\% confidence interval for the median & interval \\
$S$ & Speedup ratio (baseline / proposed) & scalar (dimensionless) \\
\bottomrule
\end{tabular}
\end{table}

\subsection{Baseline vs.\ Proposed: Execution-Time Comparison}
\label{sec:comparison}

Table~\ref{tab:main} reports median execution time, interquartile range, and speedup for the baseline and proposed aggregation methods across four input sizes.

\begin{table}[htbp]
\centering
\caption{Median execution time (ms) with interquartile range (p25--p75) and speedup for baseline and proposed aggregation methods across four input sizes. Best median per row in \textbf{bold}.}
\label{tab:main}
\begin{tabular}{rrrr}
\toprule
Input rows & Baseline median (ms) & Proposed median (ms) & Speedup $S$ \\
\midrule
10     & 3.537 (3.517--3.571)   & \textbf{1.910} (1.903--1.920) & 1.85$\times$ \\
100    & 3.635 (3.622--3.658)   & \textbf{1.929} (1.921--1.938) & 1.88$\times$ \\
1{,}000 & 5.016 (4.992--5.171) & \textbf{2.193} (2.176--2.291) & 2.29$\times$ \\
10{,}000 & 20.647 (20.481--21.025) & \textbf{4.852} (4.797--4.975) & \textbf{4.26}$\times$ \\
\bottomrule
\end{tabular}
\end{table}

Speedup is defined as $S = \mathrm{median}(x_{\text{baseline}})/\mathrm{median}(x_{\text{proposed}})$ (Equation~\eqref{eq:speedup}). A value $S > 1$ indicates that the proposed method is faster. Both methods produced matching outputs on every tested workload; all 800 measured runs completed without errors (error rate = 0\%).

\subsection{Visual Comparison}
\label{sec:visual}

The quantitative results from Table~\ref{tab:main} are visualized in Figures~\ref{fig:latency} and~\ref{fig:dist}. Figure~\ref{fig:latency} shows the median execution time as a function of input size. Figure~\ref{fig:dist} shows the full execution-time distributions for each method.

\begin{figure}[htbp]
\centering
\includegraphics[width=0.85\linewidth]{figures/fig1_latency_vs_load.pdf}
\caption{Median execution time (ms) versus input size (logarithmic scale) for the baseline and proposed aggregation methods. Vertical bars represent bootstrap 95\% confidence intervals for the median ($B = 5000$ resamples, seed = 42). The proposed method has a lower median at every tested input size; the gap widens with input size, reaching a 4.26$\times$ speedup at 10{,}000 rows. The bars indicate that the median is estimated precisely within each session, but they do not capture variation across machines or sessions.}
\label{fig:latency}
\end{figure}

\begin{figure}[htbp]
\centering
\includegraphics[width=0.85\linewidth]{figures/fig2_distribution_boxplot.pdf}
\caption{Distribution of execution time (ms) for the baseline and proposed methods across four input sizes. Each box represents the interquartile range (p25--p75); whiskers extend to the most extreme observations within 1.5$\times$IQR; points beyond the whiskers are retained as outliers. Offset markers and bars show medians and their bootstrap 95\% confidence intervals. Each panel uses its own vertical scale. The baseline exhibits a longer upper tail at 10{,}000 rows (p99 $\approx$ 33~ms), whereas the proposed method remains tightly concentrated (p99 $\approx$ 5.9~ms).}
\label{fig:dist}
\end{figure}

\subsection{Discussion and Limitations}
\label{sec:discussion}

As shown in Table~\ref{tab:main}, the proposed aggregation method consistently reduced median execution time relative to the baseline across all tested input sizes. At 10{,}000 rows, the median dropped from 20.647~ms to 4.852~ms, corresponding to a 4.26$\times$ speedup. Figure~\ref{fig:dist} further suggests that the proposed method is more stable: the baseline's 99th percentile reached approximately 33~ms, whereas the proposed method remained below 6~ms. A plausible explanation is that the baseline performs separate grouped aggregations followed by table merges, so its cost grows with the number of participant groups as well as with row count, whereas the proposed method performs a single grouped aggregation. Within the tested workload range, these results suggest that one-pass aggregation is preferable for participant-level diary aggregation at 10{,}000+ rows.

These results concern computational performance on synthetic data within a single session on one machine. They do not establish a behavioral association between short-form video use and student procrastination. The pipeline is validated only as a fast and reproducible tool for future analysis of real participant data.

\subsubsection{Threats to Validity}

\textbf{Internal validity.} All measurements were collected on a single machine in a single session. The bootstrap confidence intervals ($B = 5000$, seed = 42) describe within-session uncertainty under an independent-run assumption; they do not capture variation across machines, operating systems, or separate sessions. With $n = 100$ measurements per condition, p99 estimates are approximate and should be interpreted as indicative rather than precise. Warm-up runs ($n_{\text{warmup}} = 10$) were discarded to reduce cold-start effects, but residual order effects cannot be excluded.

\textbf{External validity.} Both methods were evaluated on synthetic rows derived from the committed 10-row sample, using a single Python/pandas implementation stack. The results therefore may not generalize to real diary data with realistic missingness, to alternative implementations (e.g., SQL-based aggregation), or to other hardware configurations. In addition, the experimental workload increased row count and participant-group count together; their separate effects were not isolated, so the observed speedup cannot be attributed to either factor alone.

\textbf{Construct validity.} Execution time was measured only for the \texttt{derive\_participant\_level} function; CSV reading, file writing, and console output were outside the timed section. The reported timings therefore reflect the aggregation step in isolation, not end-to-end pipeline latency. Speedup $S$ is a ratio of medians (Equation~\eqref{eq:speedup}) and does not by itself characterize tail behavior, which is why percentile-based summaries are reported alongside it.

\textbf{Scope of claims.} Within the tested envelope (10--10{,}000 synthetic rows, one machine, one session, pandas implementation), the results suggest that one-pass aggregation reduces median execution time by a factor of 1.85$\times$ to 4.26$\times$. These findings evaluate software performance only and should not be read as evidence for or against any behavioral hypothesis about short-form video consumption and academic procrastination.
